\section{Bisection Method}
    \subsection{Main Idea}
        \textbf{Bisection method} is based on IVT, Theorem \ref{thm:IVT}.
        Given a function $f$, we start with initial interval $[a, b]$ such that
        $f(a)f(b)<0$, checking the sign of intermediate value 
        $f(\frac{a+b}{2})$. If the sign equals that of $f(a)$, we set $a = \frac{a + b}{2}$, else we set $b = \frac{a + b}{2}$.
    
    \subsection{Pseudocode}
        \Algorithm{bisectionMethod}
        {Bisection Method}
        {
            \inputitem{$f(\cdot)$}{a function which we seek for its root} \\
            \inputitem{$a < b \in \mathbb{R}$}{endpoints, satisfiying $f(a)f(b)<0$} \\
            \inputitem{$N$}{number of iterations} \\
            \inputitem{$TOL$}{tolerance}
        }
        {
            \inputitem{$p$}{approximate solution}
        }
        {
            \For{$i = 0, 1, 2, \dots, N - 1$}
            {
                $p \leftarrow \frac{a+b}{2}$\;
                \If{$f(p)=0$ or $\frac{b - 2}{2} < TOL$}
                {
                    \Return{$p$} \tcp{procedure completed succesfully}
                }
                \ElseIf{$f(a)f(p) > 0$}
                {
                    $a \leftarrow p$\;
                }
                \Else
                {
                    $b \leftarrow p$\;
                }
            }
            \Return{failure after $N$ iterations}
        }

    \subsection{Implementation Details}
        \subsubsection{Stopping Procedure}
            One can change the stopping prodecure of algorithm \ref{alg:bisectionMethod}. With tolerance $\varepsilon > 0$,
            followings are examples of other stopping procedure:
            \begin{align}
                |p_N - p_{N-1}| &< \varepsilon \\
                \dfrac{|p_N - p_{N-1}|}{|p_N|} &< \varepsilon,\ p_N \neq 0 \\
                |f(p_N)| &< \varepsilon
            \end{align}
            But, there are some difficulties. There are sequences $\{p_n\}$ where $|p_n - p_{n-1}|$ converges to zero while $p_n\to\infty$.
            If is also possible for $f(p_n)$ to be close to zero while $p_n$ differs signifnicantly from $p$.
            The second procedure is the best when we know nothing about $f$ and $p$.

        \subsubsection{Overflow and Underflow}
            It is recommended to use following computation for calculating $\frac{a + b}{2}$:
            \begin{equation}
                p = a + \frac{b - a}{2}.
            \end{equation}
            When $b - a$ is near maximum precision, a round off error is still in presense, but it would not affect the value of $p$ so much.
            However, when $a + b$ is near maximum precision, $\frac{a+b}{2}$ might not even produce a result in interval $[a,b]$.

            Using \textbf{sign} function for determining the equality of sign can prevent overflow or underflow in the multiplication of $f(a)$ and $f(p)$.
    \subsection{Pros and Cons}
        Pros:
        \begin{itemize}
            \item Very easy to implement;
        \end{itemize}
        Cons:
        \begin{itemize}
            \item Slow convergence rate. It has linear convergence rate;
            \item Cannot extend to multi-dimensional functions;
            \item Does not exploit additional informations such as derivatives. This leads to slow convergence rates.
        \end{itemize}
    \subsection{Convergence}
            TODO: Bisection Method Convergence Proof

\section{Fixed-Point Iteration}
    \subsection{Main Idea}
        \textbf{Fixed-point iteration} is based on Theorem \ref{thm:fixedPoint}. The definition of fixed point is as following:
        \begin{defn}{Fixed Point}{fixedPoint}
            Let $X,\ Y \subset \mathbb{R}$. Then the \textbf{fixed point} of a function $f:X \to Y$ is a point $p \in X$ such that
            \begin{equation}
                f(p)=p.
            \end{equation}
        \end{defn}
        There are sufficient conditions for the existence and uniqueness of a fixed point.
        \begin{thm}{Existence and Uniqueness of a Fixed Point}{exisUniqFixedPoint}
            Let $X,\ Y \subset \mathbb{R}$. Let $a < b \in X$, and let $g:X\to Y$ be a function.
            Then, following holds:
            \begin{enumerate}
                \item If $g \in C[a,b]$ and $g(x) \in [a,b]$ for all $x \in [a,b]$, then $g$ has at least one fixed point in $[a,b]$.
                \item If, in addition, $g'(x)$ exists on $(a,b)$ and a constant $0 < k < 1$ exists with $|g'(x)| \leq k$ for all $x \in (a,b)$,
                    then there exists a unique fixed point in $(a,b)$.
            \end{enumerate}
        \end{thm}
        \begin{proof}
            Can be proved with IVT and MVT.
        \end{proof}
        One may wonder what fixed point does with the root of a function. By modifying the equation $f(x)=0$ a bit,
        we can reformulate it into an equation about fixed point, for example, $f(x) + x = x$ or $x - 3f(x) = x$.

        TODO: Figures of fixed point iteration
    \subsection{Pseudocode}
        \Algorithm{fixedPointIter}
        {Fixed-Point Iteration}
        {
            \inputitem{$g(\cdot)$}{a function which we seek for its fixed point} \\
            \inputitem{$p_0$}{initial approximation} \\
            \inputitem{$N$}{number of iterations} \\
            \inputitem{$TOL$}{tolerance}
        }
        {
            \inputitem{$p$}{approximate solution}
        }
        {
            \tcp{initialize}
            $p \leftarrow p_0$\;
            \For{$i = 0, 1, 2, \dots, N$}
            {
                $p \leftarrow g(p)$\;
                \If{$|p-p_0| < TOL$}
                {
                    \Return{$p$} \tcp{procedure completed succesfully}
                }
            }
            \Return{failure after $N$ iterations}
        }
    \subsection{Implementation Details}
    \subsection{Pros and Cons}
        Pros:
        \begin{itemize}
            \item Easy to implement;
        \end{itemize}
        Cons:
        \begin{itemize}
            \item Slow convergence rate. It has linear convergence rate. Methods of increasing convergence rate is treated in \ref{}.
            \item Depends on the modified equation. Convergence result depends on the modified equation of $f(x)=0$;
        \end{itemize}
    \subsection{Convergence}
        \subsubsection{Fixed-Point Theorem}
            \begin{thm}{Fixed-Point Theorem}{fixedPoint}
            \end{thm}